Under the two tested pressure prompts, roughly one fifth to one quarter of incorrect answers contradict stable neutral correctness; the remainder splits between inconsistent and stable-wrong answering. This is an operational decomposition, not a prevalence estimate for intentional lying. Linear detectors can rank some of these categories, but perfect cross-prompt separation is reproduced by condition identity alone and weakens substantially when both false classes receive the same pressure. A neutral-error probe preferentially scores inconsistent errors, while a task-specific probe retains moderate same-pressure signal. Evaluations should report knowledge strata, stable-wrong cases, pressure-matched correct controls, false-only contrasts, pre-answer controls, and threshold transfer before treating a ``lie detector'' score as evidence of belief--statement mismatch.
